\section{System Architecture}
\label{sec:system-architecture}

	\begin{figure}[t]
		\centering
		\includegraphics[width=0.7\columnwidth]{graphics/gen/System_Architecture_Overview_new.pdf}
		\caption{High-level architecture of the remote control.}
		\label{fig:overview}
	\end{figure}

	\begin{figure*}[!t]
		\centering
		\includegraphics[width=\textwidth]{graphics/gen/TrainClient.pdf}
		\caption{Top level components of the \textsf{TrainClient}.}
		\label{fig:train-client}
	\end{figure*}

	Figure \ref{fig:overview} illustrates the high-level structure of our system, highlighting the main modules. The \textsf{TrainClient} runs on the train, collects sensor data from the physical world and reacts upon receiving control commands. The \textsf{RemoteControlClient} is a Vue.js web application that provides a user interface for operators, rendering video streams and sending control commands. The \textsf{CentralServer} is a Python FastAPI service deployed in the cloud that acts as a communication hub between the \textsf{TrainClient} and the \textsf{RemoteControlClient}. In this paper, our focus is to design the architecture of the \textsf{TrainClient} using Lingua Franca.

	In Lingua Franca, a \emph{reactor} is a concurrent, stateful component with well-defined input and output ports, timers, and actions. Reactions are functions triggered by events, execute atomically and deterministically. Lingua Franca enforces a logical time semantics: all events are scheduled and processed in timestamp order, and physical time can be explicitly referenced for real-time interactions. Figure~\ref{fig:train-client} illustrates the top-level components of the \textsf{TrainClient}. The diagrams are automatically generated using the diagramming framework KLighD\footnote{https://github.com/kieler/KLighD} and the Eclipse Layout Kernel (ELK)~\cite{KasperowskiDvH24}.

	The \textsf{TrainClient} is the main \emph{reactor} containing all other \emph{reactors}. The \textsf{VideoStreamer} is responsible for capturing and encoding video frames from the camera. The encoded video frames are then passed to both \textsf{NetworkProtocol} and \textsf{DataStorage}. The \textsf{DataStorage} is responsible for storing data for further analysis. The \textsf{NetworkProtocol} handles the connection establishment and data transmission with the \textsf{CentralServer}. Upon data reception, the \textsf{NetworkProtocol} forwards the data to the \textsf{Parser}. When the \textsf{Parser} identifies a new connection request, it sends output signals to initialize the \textsf{ClockSynchronizer} and the \textsf{HeartbeatHandler}. The \textsf{ClockSynchronizer} estimates the clock offset between the two systems. For our purposes, a synchronization accuracy in the msec range is sufficient. For our initial design, we adopted a rather simple scheme, where the initiating side sends its local timestamp, the remote side appends its own timestamp and returns the message, and the clock offset is derived from these timestamp and the midpoint of the round trip time. This exchange is repeated several times and the resulting estimates are averaged. As future work, we plan to investigate the use of an established clock synchronization protocol, such as NTP or SNTP, or possibly GPS~\cite{SeliemRA+26}. The \textsf{HeartbeatHandler} manages the periodic exchange of heartbeat messages to monitor active connectivity. Upon receiving control commands, the \textsf{Parser} forwards them to the \textsf{Controller} which processes the commands and translates them into actionable tasks for the physical world and actuates the train accordingly.

	The \textsf{HeartbeatHandler}, shown in Figure \ref{fig:heartbeat}, demonstrates how Lingua Franca's logical-time semantics and \emph{modal reactors} support robust connectivity monitoring. Modal reactors coordinate reactive tasks based on modes of operation~\cite{SchulzRosengartenvHLB+23c}. The \textsf{HeartbeatHandler} operates in three modes: \textsf{Init}, \textsf{Send-and-Check}, and \textsf{Stopped}. Upon receiving a heartbeat initialization signal, the \emph{reactor} switches from \textsf{Init} to \textsf{Send-and-Check} and schedules a logical action with a 3000~ms delay, within which at least one heartbeat message is expected from the remote control. The LF execution model ensures that a reaction scheduled at a logical time is not executed before the corresponding physical time is reached. A 500~ms periodic timer sends heartbeat messages, and the logical action is rescheduled upon receiving a new heartbeat. If the heartbeats stop signal is received or no heartbeat arrives within the time window, a timeout signal is sent and the \emph{reactor} switches to \textsf{Stopped}, indicating loss of connectivity or inactive remote control.



	\begin{figure}[t]
		\centering
		\includegraphics[width=0.9\columnwidth]{graphics/gen/HeartbeatHandler-cropped.pdf}
		\caption{\textsf{HeartbeatHandler} monitoring connectivity through periodic heartbeat messages.}
		\label{fig:heartbeat}
	\end{figure}